\section{Spectral reductions and finite coordinate cores}
\label{sec:structure}

Throughout this section \(X=(X_1,\ldots,X_n)\) has independent
coordinates with probability densities bounded by one.  Write
\[
 Q_A(X)=X^{\mathsf T}AX+b^{\mathsf T}X+c_0,
 \qquad
 N_A(t)=\#\{j:|\lambda_j(A)|\geq t\}.
\]
All assertions are uniform in the affine terms.  For \(0<c\leq1\),
a \emph{\(c\)-pivot} is either a diagonal entry of modulus at least
\(c\), or a principal \(2\)-by-\(2\) submatrix whose determinant has
modulus at least \(c\).  If there is no \(c\)-pivot, then
\begin{equation}
 |a_{ii}|<c,\qquad |a_{ij}|^2<c+c^2\leq2c .
 \label{eq:no-pivot-entries}
\end{equation}

\subsection{Schur complements and logarithmic potentials}

\begin{lemma}[Spectral count under a coordinate pivot]
\label{lem:schur-count}
Let \(B\) be an invertible principal submatrix of a real symmetric
matrix \(A\), of order \(k\), and let \(S=A/B\) be its Schur
complement.  For every \(t>0\),
\[
 N_S(t)\geq N_A(t)-k.
\]
\end{lemma}
\begin{proof}
First suppose \(A\) is invertible.  Then \(S\) is invertible and
\(S^{-1}\) is the complementary principal submatrix of \(A^{-1}\).
Interlacing bounds the number of eigenvalues of \(S^{-1}\) in
\((t^{-1},\infty)\) by the corresponding number for \(A^{-1}\).
It gives the analogous bound in \((-\infty,-t^{-1})\).
Consequently, at most \(n-N_A(t)\) of the \(n-k\) eigenvalues of
\(S^{-1}\) have modulus greater than \(t^{-1}\).  This is the
asserted inequality.  For singular \(A\), apply the argument to
\(A+\varepsilon I\) along values for which this matrix is invertible,
at threshold \(t-|\varepsilon|\).  The pivot remains invertible
for sufficiently small \(\varepsilon\); its Schur complement
converges to \(S\).  Continuity of the ordered eigenvalues gives
the result.
\end{proof}


Define
\[
 {\cal L}(Y)=\sup_s\mathbb E\log_+\frac1{|Y-s|},
 \qquad
 {\cal H}(Y)=\sup_s\mathbb E\left[
 |Y-s|^{-1/2}\boldsymbol 1_{\{|Y-s|\leq1\}}\right].
\]

\begin{lemma}[A global logarithmic bound]
\label{lem:global-log}
If \(A\ne0\), then
\[
 {\cal L}(Q_A)\leq
 4+\log_+\frac{5^4}{\|A\|}.
\]
In particular, \(N_A(1)\geq1\) implies
\({\cal L}(Q_A)\leq L_0:=4+4\log5\).
\end{lemma}
\begin{proof}
The quarter-power estimate of Lemma~\ref{four:quarter} and the
layer-cake formula give
\[
 \mathbb E\log_+\frac1{|Q_A-s|}
 \leq\int_0^\infty
 \min\{1,5\|A\|^{-1/4}e^{-u/4}\}\,du.
\]
The integral is \(4+4\log K\) when \(K=5\|A\|^{-1/4}\geq1\),
and \(4K\leq4\) otherwise.
\end{proof}

\begin{lemma}[Half-power integrability at a pivot]
\label{lem:pivot-half}
If \(N_A(1)\geq2\) and \(A\) has a \(c\)-pivot, then
\[
 {\cal H}(Q_A)\leq44c^{-1/2}.
\]
\end{lemma}
\begin{proof}
Suppose first that \(a=a_{hh}\) satisfies \(|a|\geq c\).
Completing the square in the original coordinate \(X_h\) writes
\(Q_A=a(X_h-m)^2+K\), where \(m\) is affine in the remaining
coordinates and the quadratic matrix of \(K\) is \(A/a\).
Lemma~\ref{lem:schur-count} gives \(N_{A/a}(1)\geq1\).
The conditional density estimate of Lemma~\ref{geom:square}
therefore yields
\[
 {\cal H}(Q_A)
 \leq |a|^{-1/2}\sup_s\mathbb E J(s-K),
 \qquad
 J(d)=\int_{-1}^1 |x|^{-1/2}|x-d|^{-1/2}\,dx .
\]
For \(d\ne0\),
\begin{equation}
 J(d)\leq8+3\log_+\frac1{|d|}.
 \label{eq:root-convolution}
\end{equation}
Indeed, on \(|x|\leq2|d|\) substitute \(x=|d|u\).
The resulting integral is at most
\[
 \int_{-2}^2 |u|^{-1/2}|u-1|^{-1/2}\,du
 =2\operatorname{arsinh}\sqrt2+\pi
   +2\operatorname{arcosh}\sqrt2<8.
\]
On \(2|d|<|x|\leq1\), the integrand is at most
\(\sqrt2/|x|\), giving the remaining term in
\eqref{eq:root-convolution}.  Lemma~\ref{lem:global-log}
bounds this case by \(c^{-1/2}(8+3L_0)<40c^{-1/2}\).

Next suppose that a principal two-coordinate block \(B\) is
definite with determinant \(D\geq c\).  Its conditional density
is at most \(\pi D^{-1/2}\); integrating
\(|x-s|^{-1/2}\boldsymbol1_{\{|x-s|\leq1\}}\) gives
\({\cal H}(Q_A)\leq4\pi c^{-1/2}\).

Finally suppose \(\det B=-D\), \(D\geq c\).
Lemma~\ref{geom:pair} bounds the conditional density by
\[
 D^{-1/2}\left(5.2+
       \log_+\frac{16\sqrt D}{|y-K|}\right).
\]
For every real \(a\),
\[
 \int_{-1}^1 |x|^{-1/2}\log_+\frac1{|x-a|}\,dx\leq8.
\]
To see this without a regularity assumption, apply layer cake
to both factors.  Their superlevel sets are intervals, and the
overlap of two intervals of prescribed lengths is maximal when
their centres coincide.  The integral is therefore at most its
value at \(a=0\), namely
\(2\int_0^1x^{-1/2}\log(1/x)\,dx=8\).
It follows that
\[
 {\cal H}(Q_A)\leq
 D^{-1/2}\left(28.8+4\log_+(16\sqrt D)\right)
 \leq c^{-1/2}(28.8+4\log16)+4/e
 <44c^{-1/2}.
\]
All conditionings are on original independent coordinates.
\end{proof}

\subsection{Gaussian padding and reduction to small pivots}

\begin{theorem}[Reduction to the three-eigenvalue no-pivot case]
\label{thm:no-pivot-equivalence}
Fix \(0<c\leq1\).  Suppose that
\begin{equation}
 \|p_{Q_A}\|_\infty\leq C_3
 \quad\text{whenever }N_A(1)\geq3
 \text{ and \(A\) has no \(c\)-pivot}.
 \label{eq:no-pivot-hypothesis}
\end{equation}
Then the bounded-density conjecture holds for all real symmetric
matrices with \(N_A(1)\geq3\).  One admissible bound is
\[
 \max\left\{C_3,\ c^{-1/2}(1+G_2),\ 20c^{-1/2}\right\},
 \qquad
 G_2=\max\{2\sqrt{2\pi e}\,C_3,\ 44c^{-1/2}\}.
\]
The converse implication is immediate.
\end{theorem}
\begin{proof}
First consider a matrix \(M\) with \(N_M(1)\geq2\) and no
\(c\)-pivot.  Choose \(N>1/c\), put
\(u=N^{-1/2}(1,\ldots,1)\), and let \(G_1,\ldots,G_N\) be
independent standard Gaussian variables, independent of \(X\).
The variable \(Z=\langle u,G\rangle\) is standard Gaussian.
The matrices \(M\oplus uu^{\mathsf T}\) and
\(M\oplus(-uu^{\mathsf T})\) have at least three eigenvalues
of modulus at least one and no \(c\)-pivot: all new diagonal
entries have modulus \(1/N<c\), all two-coordinate minors
inside the new block vanish, and a mixed minor is the product
of two diagonal entries of modulus less than \(c\).
The original Gaussian coordinate densities are bounded by one.
Thus \eqref{eq:no-pivot-hypothesis} applies to \(Q_M\pm Z^2\).

The density of \(Z^2\) is
\[
 k(t)=\frac{e^{-t/2}}{\sqrt{2\pi t}}\boldsymbol1_{\{t>0\}}.
\]
For \(0<|t|\leq1\),
\[
 k(t)+k(-t)\geq(2\pi e)^{-1/2}|t|^{-1/2}.
\]
Adding the two convolution bounds gives
\({\cal H}(Q_M)\leq2\sqrt{2\pi e}\,C_3\).
The estimate holds for every centre \(s\), rather than only
almost every centre: the symmetrized kernel
\((2\pi|t|)^{-1/2}e^{-|t|/2}\), with value \(+\infty\) at zero,
is lower semicontinuous.  Its convolution has the same property,
so an almost-everywhere upper bound extends everywhere.
Lemma~\ref{four:quarter} also ensures that \(Q_M\) has no atoms.
Together with Lemma~\ref{lem:pivot-half}, this proves
\begin{equation}
 N_M(1)\geq2\quad\Longrightarrow\quad {\cal H}(Q_M)\leq G_2
 \label{eq:global-half-conditional}
\end{equation}
for every matrix \(M\).

Now let \(N_A(1)\geq3\).  The case with no pivot is the
hypothesis.  At a one-coordinate pivot \(a\), the square
kernel and Lemma~\ref{lem:schur-count} give
\[
 p_{Q_A}(s)\leq |a|^{-1/2}\mathbb E|s-K|^{-1/2}
 \leq c^{-1/2}(1+G_2),
\]
using \eqref{eq:global-half-conditional} for the Schur
complement.  A definite two-coordinate pivot gives
\(\pi c^{-1/2}\).  At an indefinite pivot of determinant
\(-D\), the centre has at least one large eigenvalue, so
Lemmas~\ref{geom:pair} and \ref{lem:global-log} give
\[
 p_{Q_A}(s)\leq
 D^{-1/2}\bigl(5.2+\log_+(16\sqrt D)+L_0\bigr)
 \leq20c^{-1/2}.
\]
This exhausts the cases.
\end{proof}

\begin{remark}
The reduction assumes \eqref{eq:no-pivot-hypothesis} and requires
neither regularity of the original coordinate densities nor moment
bounds.
\end{remark}

\subsection{Finite-dimensional bounds and projected gradients}

\begin{theorem}[Slicing by a projected gradient]
\label{thm:projected-gradient-slicing}
Let \(E\) be a three-dimensional invariant subspace of \(A\)
on which every eigenvalue has modulus at least \(\sigma>0\).
Put
\[
 \ell_i=|P_Ee_i|^2,\qquad {\cal R}(E)=\sum_i\sqrt{\ell_i}.
\]
Then
\begin{equation}
 \|p_{Q_A}\|_\infty
 \leq\frac{2\sqrt{2\pi}}{\sigma}{\cal R}(E)
 \leq\frac{2\sqrt{6\pi n}}{\sigma}.
 \label{eq:projected-slicing}
\end{equation}
For \(n=3\), the bound improves to \(6\sqrt\pi/\sigma\).
All eigenvalues outside \(E\) are unrestricted.
\end{theorem}
\begin{proof}
Write \(v(x)=\nabla Q_A(x)\) and \(g(x)=P_Ev(x)\).
There is a unique \(z_0\in E\) with
\(2(A|_E)z_0+P_Eb=0\), and hence
\[
 g(x)=2(A|_E)(P_Ex-z_0),\qquad
 |g(x)|\geq2\sigma|P_Ex-z_0|.
\]
The set \(g=0\) has codimension three and has probability zero.
The identities \(v^{\mathsf T}g=|g|^2\) and \(g\in E\) imply
\begin{equation}
 1\leq\sum_i\frac{|\partial_iQ_A|\,|g_i|}{|g|^2},
 \qquad |g_i|\leq\sqrt{\ell_i}|g|.
 \label{eq:projected-field-partition}
\end{equation}
The absolute values in this inequality are necessary; individual
terms of \(v^{\mathsf T}g\) need not be nonnegative.

Let \(\psi\) be a nonnegative bounded compactly supported
Borel function.  Multiply \eqref{eq:projected-field-partition}
by \(\psi(Q_A(x))\prod_jf_j(x_j)\), integrate, and apply
Tonelli.  For the \(i\)-th summand, fix all coordinates except
\(x_i=t\).  On the monotonicity intervals of the resulting
nonconstant quadratic or linear polynomial, one-dimensional
change of variables gives
\[
 \int \psi(Q_A)f_i(t)
       \frac{|\partial_iQ_A|\,|g_i|}{|g|^2}\,dt
 =
 \int\psi(s)
 \sum_{\substack{t:Q_A(t,x_{-i})=s\\ \partial_iQ_A\ne0}}
 f_i(t)\frac{|g_i(t,x_{-i})|}{|g(t,x_{-i})|^2}\,ds .
\]
Critical values form a null set.  If the fibre polynomial is
constant, the left side is zero and the fibre contributes zero.

Choose a two-dimensional subspace \(E_i\subset E\cap e_i^\perp\).
It exists because \(\dim E=3\).  The projection
\(P_{E_i}x\) is independent of the coordinate \(t=x_i\), and
\[
 |g(x)|\geq2\sigma|P_{E_i}x-P_{E_i}z_0|.
\]
There are at most two noncritical roots, and \(f_i\leq1\).
The root sum is consequently at most
\[
 \frac{\sqrt{\ell_i}}
 {\sigma|P_{E_i}x_{-i}-P_{E_i}z_0|}.
\]
The vector \(P_{E_i}X\) is an orthonormal two-dimensional
projection of the remaining original product coordinates.
Its density is at most \(2\), by
\citet[Theorem~1.1]{LPP2016}.
For a two-dimensional probability density bounded by \(M\),
layer cake gives
\[
 \mathbb E|Y-z|^{-1}
 \leq\int_0^\infty\min\{1,\pi M/u^2\}\,du
 =2\sqrt{\pi M}.
\]
The \(i\)-th summand is therefore bounded by
\((2\sqrt{2\pi}/\sigma)\sqrt{\ell_i}\int\psi\).
Summing proves domination of the output law by the first
constant in \eqref{eq:projected-slicing} times Lebesgue
measure.  This establishes absolute continuity as well as
the essential density bound, without a pointwise density
assumption on a level set.

Finally \(\sum_i\ell_i=3\), so
\({\cal R}(E)\leq\sqrt{3n}\).
When \(n=3\), \(E_i\) is the full original coordinate
complement, whose density is at most one; also
\({\cal R}(E)=3\).  This gives \(6\sqrt\pi/\sigma\).
\end{proof}

\begin{corollary}[The finite-dimensional spectral problem]
\label{thm:finite-dimension}
If \(A\) is \(n\)-by-\(n\) and \(N_A(\lambda)\geq3\), then
\[
 \|p_{Q_A}\|_\infty\leq2\sqrt{6\pi n}/\lambda.
\]
\end{corollary}
\begin{proof}
Choose \(E\) to be spanned by three corresponding eigenvectors
and apply Theorem~\ref{thm:projected-gradient-slicing}.
\end{proof}

\begin{corollary}[A finite coordinate core]
\label{cor:finite-core}
Suppose a principal submatrix \(A_{KK}\), with \(|K|\leq N\),
has three eigenvalues of modulus at least \(\lambda>0\).
Then
\[
 \|p_{Q_A}\|_\infty\leq2\sqrt{6\pi N}/\lambda,
\]
irrespective of the ambient dimension and all other blocks of \(A\).
\end{corollary}
\begin{proof}
Condition on \(X_{K^c}\).  The remaining variables retain their
independent densities, the quadratic matrix is \(A_{KK}\),
and the conditioned values affect only its affine terms.
Apply Corollary~\ref{thm:finite-dimension} and average.
\end{proof}

\begin{theorem}[A localized spectral core with a dense perturbation]
\label{thm:physical-core-stability}
Let \(H\) contain \(m\geq3\) original coordinates, and suppose
\[
 A_{HH}=U_H\Lambda U_H^{\mathsf T}+T_H,\qquad
 U^{\mathsf T}U=I_3,\qquad
 U_H^{\mathsf T}U_H\geq aI_3,
\]
where \(\Lambda\) is symmetric,
\(s_{\min}(\Lambda)\geq s>0\), and
\(\|T_H\|\leq\varepsilon<as\).
Then every choice of the other blocks of \(A\) satisfies
\[
 \|p_{Q_A}\|_\infty\leq
 \frac{2\sqrt{6\pi m}}{as-\varepsilon}.
\]
\end{theorem}
\begin{proof}
Put \(G=U_H^{\mathsf T}U_H\) and
\(V=U_HG^{-1/2}\).  The columns of \(V\) are orthonormal.
On their span, the rank-three matrix
\(U_H\Lambda U_H^{\mathsf T}\) is represented by
\(K=G^{1/2}\Lambda G^{1/2}\), for which
\[
 |Kz|\geq\sqrt a\,|\Lambda G^{1/2}z|
 \geq as|z|.
\]
For \(x\in\operatorname{range}(V)\), the triangle inequality
gives \(|A_{HH}x|\geq(as-\varepsilon)|x|\).
The variational characterization of singular values yields
\(s_3(A_{HH})\geq as-\varepsilon\).
Now apply Corollary~\ref{cor:finite-core}.
\end{proof}

The perturbation in Theorem~\ref{thm:physical-core-stability}
is retained in the actual physical matrix.  Its small operator
norm certifies the spectrum on the core; the proof makes no
assertion that \(\varepsilon|X|^2\) is small.

\begin{proposition}[Small balls of a projected gradient]
\label{prop:projected-gradient}
Let \(u_1,\ldots,u_r\) be orthonormal eigenvectors of \(A\),
with nonzero eigenvalues \(\lambda_1,\ldots,\lambda_r\),
and put \(U=(u_1,\ldots,u_r)\).
The vector \(Y=U^{\mathsf T}\nabla Q_A(X)\) has density at most
\[
 K=\frac{2^{-r/2}}{\prod_{j=1}^r|\lambda_j|}.
\]
Writing \(v_r\) for the volume of the unit ball in
\(\mathbb R^r\), for every \(z\in\mathbb R^r\), \(a>0\),
and \(0<p<r\),
\[
 \mathbb P\{|Y-z|\leq a\}\leq\min\{1,Kv_ra^r\},
 \qquad
 \mathbb E|Y-z|^{-p}
 \leq\frac r{r-p}(Kv_r)^{p/r}.
\]
\end{proposition}
\begin{proof}
Since
\[
 U^{\mathsf T}\nabla Q_A
 =2\operatorname{diag}(\lambda_1,\ldots,\lambda_r)U^{\mathsf T}X
   +U^{\mathsf T}b,
\]
the projection-density bound of Lemma~\ref{four:projection}
and change of variables give \(K\).
Integrating this density over a ball gives the first estimate.
The second follows by integrating the first estimate in the
layer-cake identity
\(\mathbb E|Y-z|^{-p}
=\int_0^\infty\mathbb P\{|Y-z|<t^{-1/p}\}\,dt\).
\end{proof}

The projected-gradient estimate controls neighbourhoods of
critical points.  It does not, by itself, bound the density of
\(Q_A\): the level-set integral also contains the product
density, which is allowed to be rough.

\subsection{An affine estimate for exact rank two}

\begin{lemma}[A square with an arbitrary linear term]
\label{lem:rank-one-affine}
For \(|u|=1\), \(\mu\ne0\), and arbitrary \(b,c_0\),
\[
 \sup_s\mathbb P\{|\mu\langle u,X\rangle^2+
 b^{\mathsf T}X+c_0-s|\leq a\}
 \leq8\sqrt{a/|\mu|}.
\]
\end{lemma}
\begin{proof}
This is Lemma~\ref{rt:rank-one}; its Gaussian smoothing
argument includes arbitrary original-coordinate affine terms.
The closed-interval version follows by decreasing the radius
to \(a\) from above.
\end{proof}

\begin{theorem}[Exact rank two]
\label{thm:rank-two-affine}
Suppose \(A=RMR^{\mathsf T}\), where \(R^{\mathsf T}R=I_2\),
\(M\) is real symmetric, and \(|\det M|\geq1\).
Then, for \(0<a\leq1\),
\[
 \sup_s\mathbb P\{|Q_A-s|\leq a\}
 \leq32\sqrt2\,a(1+\log(1/a))\leq46a(1+\log(1/a)).
\]
Moreover, Theorem~\ref{rt:half} gives the stronger fractional-potential bound
\(\sup_s\mathbb E|Q_A-s|^{-1/2}<68\).
\end{theorem}
\begin{proof}
Let \(\rho_i=R^{\mathsf T}e_i\).
If every \(|\rho_i|^2\leq1/2\), Fresnel linearization in the
two spectral directions and the geometric
Brascamp--Lieb estimate give
\(|\phi_{Q_A}(t)|\leq4\pi/(|t|\sqrt{|\det M|})\).
The Fej\'er kernel
\((\sin(x/2)/(x/2))^2\) is greater than \(9/10\) on \([-1,1]\).
Its triangular Fourier representation yields
\[
 \mathbb P\{|Y-s|\leq a\}
 \leq\frac{10}{9}a\int_{|t|\leq a^{-1}}|\phi_Y(t)|\,dt.
\]
Using \(|\phi|\leq1\) on \(|t|\leq1\) gives the asserted
bound with constant less than \(28\).

Otherwise choose \(h\) with \(l=|\rho_h|^2>1/2\).
Rotate the two-dimensional frame so that
\(\rho_h=\sqrt l\,e_2\), and write
\[
 M=\begin{pmatrix}s&q\\q&a_0\end{pmatrix}.
\]
The first column \(r_1\) of \(R\), restricted to \(h^c\), is
a unit vector; the second restricted column \(r_2\) is
orthogonal to it.  As a quadratic in \(x=X_h\), write
\(Q_A=ax^2+2Lx+T\), and define
\(\Delta=L^2+a(s_0-T)\), for a target level \(s_0\).
Its quadratic matrix is
\[
 cc^{\mathsf T}-aA_{h^c,h^c}
 =-l\det(M)\,r_1r_1^{\mathsf T},
 \qquad
 c=A_{h^c,h}=\sqrt l(qr_1+a_0r_2),\quad a=la_0.
\]
This identity follows by expanding both sides; the coefficients
of \(r_2r_2^{\mathsf T}\) and
\(r_1r_2^{\mathsf T}+r_2r_1^{\mathsf T}\) cancel.
The other terms of \(\Delta\) are affine.
Lemma~\ref{lem:rank-one-affine} thus gives
\(\mathbb P\{|\Delta-t|\leq u\}\leq8\sqrt2\sqrt u\).

The one-variable quadratic sublevel estimate, including the
linear case \(a=0\), bounds the conditional probability by
\(\min\{1,4a_{\!s}/\sqrt{|\Delta|}\}\), where \(a_{\!s}\)
is the small-ball radius.  Hence
\[
 \mathbb P\{|Q_A-s_0|\leq a_{\!s}\}
 \leq\int_0^1\min\{1,32\sqrt2\,a_{\!s}/\theta\}\,d\theta
 \leq32\sqrt2\,a_{\!s}(1+\log(1/a_{\!s})).
\]
The potential estimate is Theorem~\ref{rt:half}, since
\(|\det M|\geq1\).
\end{proof}

\section{Off-diagonal rank and quantitative skeletons}
\label{sec:skeletons}

For a real symmetric matrix \(A\), let
\[
 \rho(A)=\max_{S\cap T=\varnothing}\operatorname{rank}A_{ST},
 \qquad
 \operatorname{mr}(A)=\min_{D\ {\rm diagonal}}\operatorname{rank}(A-D).
\]
The sets \(S,T\) in this definition need not have the same size.

\begin{theorem}[Exact completion by a diagonal]
\label{thm:diagonal-completion}
For every real symmetric \(A\),
\[
 \operatorname{mr}(A)\leq3\rho(A).
\]
For every positive integer \(r\), examples satisfy
\(\rho(A)=r\) and \(\operatorname{mr}(A)=2r\).
If \(\rho(A)\leq1\), then \(\operatorname{mr}(A)\leq2\).
\end{theorem}
\begin{proof}
Put \(r=\rho(A)\).  The case \(r=0\) means that \(A\) is
diagonal.  Otherwise choose disjoint \(r\)-sets \(S,T\)
with \(A_{ST}\) invertible, and put \(C=A_{ST}^{-1}\),
\(K=S\cup T\), \(L=K^c\).
For \(l\in L\), write \(x_l=A_{lT}\), \(y_l=A_{lS}\),
and collect these rows in matrices \(X,Y\).
Every off-diagonal \((r+1)\)-minor vanishes; adjoining
distinct \(l,m\in L\) to the selected minor gives
\[
 a_{lm}=x_lCy_m^{\mathsf T}.
\]
Symmetry therefore implies that the span of the vectors
\((y_l,x_l)\in\mathbb R^{2r}\) is isotropic for the
nondegenerate alternating form
\[
 \omega((y,x),(y',x'))=xC(y')^{\mathsf T}-x'Cy^{\mathsf T}.
\]
An isotropic subspace has dimension at most \(r\):
it is contained in its orthogonal complement for this
nondegenerate form, and the two dimensions sum to \(2r\).
The matrix \(P=XCY^{\mathsf T}\) is symmetric because its
off-diagonal entries agree with those of \(A_{LL}\).
Choose \(D_K=0\) and
\(D_L=\operatorname{diag}(A_{LL}-P)\).
Each \(L\)-row of \(A-D\) is a fixed linear image of
\((y_l,x_l)\), so these rows span at most \(r\) dimensions.
The \(2r\) rows in \(K\) add at most \(2r\), proving
\(\operatorname{mr}(A)\leq3r\).

A star with centre \(h\), zero leaf--leaf entries and at
least two nonzero edges has \(\rho=1\) and
\(\operatorname{mr}=2\).  Its off-diagonal matrix already
has rank two.  A symmetric rank-one completion would have
the form \(\pm vv^{\mathsf T}\); two nonzero centre--leaf
entries would then force a nonzero entry between those
leaves.  Direct sums of \(r\) such stars have the stated
parameters, since both quantities add over direct sums.

It remains to prove the improved upper bound for \(\rho\leq1\).
Consider the support graph of the off-diagonal entries.
For four distinct indices the hypothesis reads
\(a_{ik}a_{jl}=a_{il}a_{jk}\).
Two disjoint edges consequently force all six edges on their
four endpoints.  A vertex adjacent to a clique of order at
least three is adjacent to every vertex of that clique:
pair its incident edge with edges between the other clique
vertices and use the same identity.
Thus, if disjoint edges occur, all nonzero edges lie in a
single complete graph with at least four vertices.
On this graph the numbers
\[
 c_i=\frac{a_{ij}a_{ik}}{a_{jk}}
\]
are independent of the distinct choices \(j,k\ne i\).
The four-index identities give \(c_ic_j=a_{ij}^2>0\);
all \(c_i\) have a common sign.  The edge signs obey the
same identities, so choices of signs for
\(v_i=\pm\sqrt{|c_i|}\) give \(a_{ij}=\pm v_iv_j\)
with a common overall sign.  This supplies a rank-one
diagonal completion.  If there are no disjoint edges,
the graph is a star, a triangle, or a subgraph of one.
A star has the rank-two completion already described.
For a triangle, subtracting any eigenvalue from the three
diagonal entries makes the corresponding block singular,
giving rank at most two.
\end{proof}

For \(r\geq1\), define
\[
 V_r(A)=
 \sum_{\substack{|S|=|T|=r\\S\cap T=\varnothing}}
 \det(A_{ST})^2,\qquad V_0(A)=1.
\]
The sum is over ordered pairs of subsets; each subset is
listed in its natural order.

\begin{lemma}[Volume-sampled skeleton inequality]
\label{lem:skeleton}
Suppose \(V_r(A)>0\).  Sample disjoint \(r\)-sets \(S,T\)
with probability \(\det(A_{ST})^2/V_r(A)\), omitting zero
weights, and put
\[
 P^{ST}=A_{\cdot T}A_{ST}^{-1}A_{S\cdot}.
\]
Then
\begin{equation}
 \mathbb E_{S,T}
 \sum_{\substack{l,m\notin S\cup T\\l\ne m}}
 (A-P^{ST})_{lm}^2
 \leq(r+1)^2\frac{V_{r+1}(A)}{V_r(A)}.
 \label{eq:skeleton-inequality}
\end{equation}
\end{lemma}
\begin{proof}
For every sampled nonsingular seed and every allowed \(l,m\),
the bordered determinant formula gives
\[
 \det(A_{S+l,T+m})^2
 =\det(A_{ST})^2(A-P^{ST})_{lm}^2.
\]
Sum over the seeds and \(l,m\).  A disjoint pair of
\((r+1)\)-sets is counted at most \((r+1)^2\) times,
according to which row and column are distinguished.
Some of these deletions may leave a singular seed and are
then absent from the sum.  Division by \(V_r(A)\) proves
\eqref{eq:skeleton-inequality}.
\end{proof}

\begin{remark}[Why equality is not automatic]
\label{rem:skeleton-zero-seeds}
Take
\[
 A=\begin{pmatrix}
 0&0&1&0\\0&0&0&1\\1&0&0&0\\0&1&0&0
 \end{pmatrix},\qquad r=1.
\]
Then \(V_1=V_2=4\).  Every nonzero seed has squared
off-diagonal residual sum \(2\), so the left side of
\eqref{eq:skeleton-inequality} is \(2\), whereas the
right side is \(4\).  Singular seeds may have nonsingular
extensions.  Thus the valid general statement is an
inequality, which suffices for the following consequence.
\end{remark}

\begin{theorem}[A robust low-rank completion]
\label{thm:robust-skeleton}
Suppose \(V_r(A)>0\) and
\(V_{r+1}(A)\leq\kappa V_r(A)\).
There is a decomposition
\[
 A=D+R+E,\qquad
 \operatorname{rank}R\leq4r,\qquad
 \|E\|_{\mathrm F}\leq3(r+1)\sqrt\kappa,
\]
where \(D\) is diagonal and \(E\) is symmetric, has zero
diagonal, and vanishes on \(2r\) selected rows and columns.
\end{theorem}
\begin{proof}
Choose a seed from Lemma~\ref{lem:skeleton} for which the
off-diagonal residual on \(L=(S\cup T)^c\) has squared
Frobenius norm at most \((r+1)^2\kappa\).
Write
\[
 P=A_{LT}A_{ST}^{-1}A_{SL},\qquad
 F=\operatorname{offdiag}(A_{LL}-P).
\]
Since \(A_{LL}\) is symmetric,
\(\|P-P^{\mathsf T}\|_{\mathrm F}\leq2\|F\|_{\mathrm F}\).
Thus \(P_{\rm sym}=(P+P^{\mathsf T})/2\) satisfies
\(\|P_{\rm sym}-P\|_{\mathrm F}\leq\|F\|_{\mathrm F}\).
Let \(P_r\) retain the \(r\) largest eigenvalues in modulus
of \(P_{\rm sym}\).  The best rank-\(r\) approximation
property and \(\operatorname{rank}P\leq r\) imply
\[
 \|P_{\rm sym}-P_r\|_{\mathrm F}\leq\|F\|_{\mathrm F}.
\]
All columns of \(P_r\) lie in
\(\operatorname{col}[A_{LS},A_{LT}]\), a space of
dimension at most \(2r\), because this contains the range
of \(P_{\rm sym}\).

With \(K=S\cup T\), set
\[
 R=\begin{pmatrix}A_{KK}&A_{KL}\\A_{LK}&P_r\end{pmatrix},
 \quad
 E_{LL}=\operatorname{offdiag}(A_{LL}-P_r),
 \quad E_{K,\cdot}=E_{\cdot,K}=0,
\]
and let \(D=A-R-E\).  The \(L\)-row block of \(R\) has
rank at most \(2r\), and its \(K\)-row block adds at most
\(2r\), so \(\operatorname{rank}R\leq4r\).
The matrix \(D\) is diagonal and
\[
 \|E\|_{\mathrm F}
 \leq\|F\|_{\mathrm F}+\|P-P_r\|_{\mathrm F}
 \leq3\|F\|_{\mathrm F}
 \leq3(r+1)\sqrt\kappa.
\]
\end{proof}

\begin{corollary}
\label{cor:tail-skeleton}
If every coordinate partition \(I,J\) satisfies
\(\sum_{j\geq r+1}s_j(A_{IJ})^2\leq\delta^2\), and
\(V_r(A)>0\), then Theorem~\ref{thm:robust-skeleton}
applies with \(\kappa=4\delta^2\).  In particular
\(\|E\|_{\mathrm F}\leq6(r+1)\delta\).
\end{corollary}
\begin{proof}
For nonincreasing nonnegative numbers \(x_j\),
\[
 e_{r+1}(x)\leq
 \left(\sum_{j\geq r+1}x_j\right)e_r(x).
\]
Assign each \((r+1)\)-set to its largest index to verify
this inequality term by term.
For a random fair coordinate partition, Cauchy--Binet gives
\(\mathbb Ee_r(A_{IJ}A_{IJ}^{\mathsf T})=4^{-r}V_r(A)\).
Average the preceding inequality with \(x_j=s_j(A_{IJ})^2\).
\end{proof}

The conclusion of Corollary~\ref{cor:tail-skeleton} uses a
bound on the entire squared singular-value tail.  A bound
only on \(s_{r+1}\) yields a dimension factor through that
tail and is not the same hypothesis.

\section{Grouped projections and projected crossing blocks}
\label{sec:grouped-crossing}

The following arguments retain arbitrary phases inside each
coordinate group.  The projected coordinates need not be
independent.  This distinction is useful when many original
coordinates share a small number of coupling directions.

\begin{lemma}[A phase-weighted projection]
\label{lem:phase-projection}
Let a random vector \(W\) have density \(\rho\leq M\), and
let \(Z\) be any random variable on the same probability
space with \(|Z|\leq1\).  The complex measure
\(\mathbb E[Z;W\in dw]\) has a density \(a(w)\) satisfying
\[
 |a(w)|\leq\rho(w),\qquad
 \|a\|_1\leq1,\qquad \|a\|_2^2\leq M.
\]
\end{lemma}
\begin{proof}
Disintegration gives
\(a(w)=\rho(w)\mathbb E[Z\mid W=w]\) almost everywhere.
Its modulus is at most \(\rho\), and
\(\int|a|^2\leq\int\rho^2\leq M\int\rho=M\).
\end{proof}

\begin{theorem}[An exact-rank crossing block]
\label{thm:exact-crossblock}
Let \(X\) and \(Y\) be independent random vectors and
\(B=U\Sigma V^{\mathsf T}\) have exact rank \(r\), where
\(U,V\) have orthonormal columns and
\(\Sigma=\operatorname{diag}(s_1,\ldots,s_r)>0\).
Suppose \(U^{\mathsf T}X\) and \(V^{\mathsf T}Y\) have
densities bounded by \(M_X,M_Y\).
For arbitrary measurable real functions \(F,G\), put
\[
 Q=F(X)+2X^{\mathsf T}BY+G(Y).
\]
Then, for \(t\ne0\),
\[
 |\phi_Q(t)|\leq
 \min\left\{1,
 \left(\frac{\pi}{|t|}\right)^{r/2}
 \frac{\sqrt{M_XM_Y}}{\sqrt{s_1\cdots s_r}}\right\}.
\]
If \(r>2\), its density satisfies
\[
 \|p_Q\|_\infty\leq
 \frac r{r-2}
 \left(\frac{M_XM_Y}{s_1\cdots s_r}\right)^{1/r}.
\]
For product inputs with coordinate densities bounded by one,
one may take \(M_X,M_Y\leq2^{r/2}\).
\end{theorem}
\begin{proof}
Push \(e^{itF(X)}d\mathbb P_X\) onto \(U^{\mathsf T}X\),
obtaining a complex density \(a_t\), and similarly obtain
\(b_t\) from \(e^{itG(Y)}d\mathbb P_Y\).
Lemma~\ref{lem:phase-projection} gives
\(\|a_t\|_2^2\leq M_X\) and
\(\|b_t\|_2^2\leq M_Y\).
Independence and Fubini give
\[
 \phi_Q(t)=\int_{\mathbb R^r}
       \widehat a_t(2t\Sigma z)b_t(z)\,dz .
\]
By Plancherel,
\[
 \int|\widehat a_t(2t\Sigma z)|^2\,dz
 =\left(\frac{\pi}{|t|}\right)^r
   (s_1\cdots s_r)^{-1}\|a_t\|_2^2 .
\]
Cauchy--Schwarz proves the characteristic-function bound.
For \(\alpha>1\),
\[
 \frac1{2\pi}\int_{\mathbb R}\min\{1,K|t|^{-\alpha}\}\,dt
 =\frac{\alpha}{\pi(\alpha-1)}K^{1/\alpha}.
\]
Take \(\alpha=r/2\).  The last assertion is the projection
bound of \citet{LPP2016}.
\end{proof}

We next allow the light crossing matrix to have arbitrary
rank.  For a matrix \(C\) with rows \(c_i\), let
\[
 Z_\ell(C)=e_\ell(CC^{\mathsf T})
 =\sum_{|S|=\ell}\det(C_SC_S^{\mathsf T}).
\]
When \(Z_\ell(C)>0\), its \(\ell\)-volume row marginals are
\[
 \alpha_i=\sum_{S\ni i}
 \frac{\det(C_SC_S^{\mathsf T})}{Z_\ell(C)}.
\]
They satisfy \(\sum_i\alpha_i=\ell\).

\begin{lemma}[A weighted-subset Fourier estimate]
\label{lem:weighted-subset}
Let \(D\in\mathbb R^m\) have independent coordinate
densities bounded by one.  Let \(M\) be a \(k\)-by-\(m\)
matrix of row rank \(k\), with product of singular values
\(S_M\).  Write \(\Pi\) for projection onto its row space.
For a matrix \(A_0\) with rows \(a_i\), put
\(C=A_0(I-\Pi)\).
Suppose \(Z_\ell(C)>0\) and every \(\ell\)-volume row
marginal is at most \(1/2\).
Let \(\psi_i\) be characteristic functions of scalar
probability densities bounded by one, and let \(\kappa\)
be any nonnegative probability density on \(\mathbb R^k\).
Then
\begin{equation}
 \mathbb E\left[\kappa(MD)
     \prod_i|\psi_i(2t\langle a_i,D\rangle)|\right]
 \leq
 \frac{2^{(k+\ell)/2}(2\pi)^\ell e^{\ell/2}}
 {S_M\sqrt{Z_\ell(C)}\,|t|^\ell}.
 \label{eq:weighted-subset}
\end{equation}
When \(\ell=1\), the factor \(e^{\ell/2}\) may be omitted.
\end{lemma}
\begin{proof}
Set
\[
 w_S=\frac{\det(C_SC_S^{\mathsf T})}{Z_\ell(C)}.
\]
Discard factors with \(\alpha_i=0\), which only increases
the left side.  Weighted H\"older for the finite measure
\(\kappa(Md)d\mathbb P_D(d)\) bounds the left side by
\(\prod_{S:w_S>0}I_S^{w_S}\), where
\[
 I_S=\mathbb E\left[\kappa(MD)
       \prod_{i\in S}
       |\psi_i(2t\langle a_i,D\rangle)|^{1/\alpha_i}\right].
\]
Indeed, the exponent of each factor after taking this
weighted product is
\(\sum_{S\ni i}w_S/\alpha_i=1\).

For \(i\in S\), write \(a_i=c_i+\lambda_iM\).
The linear image
\((MD,(\langle c_i,D\rangle)_{i\in S})\) has density
bounded by
\[
 \frac{2^{(k+\ell)/2}}
 {S_M\sqrt{\det(C_SC_S^{\mathsf T})}}.
\]
To verify the determinant, the rows of \(C\) are
perpendicular to the rows of \(M\), so the Gram determinant
of the combined map is
\(S_M^2\det(C_SC_S^{\mathsf T})\).
Factor this map through an orthonormal projection and
apply Lemma~\ref{four:projection}.
For each fixed \(x=MD\), shift the remaining coordinates
by \(\lambda_i x\), integrate \(\kappa\), and use
Lemma~\ref{four:lp} at exponent \(1/\alpha_i\geq2\).
The result is
\[
 I_S\leq
 \frac{2^{(k+\ell)/2}(2\pi)^\ell}
 {S_M|t|^\ell}
 \frac{\prod_{i\in S}\sqrt{\alpha_i}}
 {\sqrt{\det(C_SC_S^{\mathsf T})}}.
\]
The weighted product of the last factors is
\[
 Z_\ell(C)^{-1/2}
 \exp\left\{\frac12\left(
       H(w)+\sum_i\alpha_i\log\alpha_i\right)\right\},
\]
where \(H(w)=-\sum_Sw_S\log w_S\).
Entropy subadditivity for the indicator coordinates of
the random subset \(S\) gives
\[
 H(w)\leq\sum_i\left[-\alpha_i\log\alpha_i
 -(1-\alpha_i)\log(1-\alpha_i)\right].
\]
This also follows directly from nonnegativity of relative
entropy against the independent Bernoulli distribution
with parameters \(\alpha_i\).
Since \(-(1-a)\log(1-a)\leq a\) and
\(\sum_i\alpha_i=\ell\), the exponential is at most
\(e^{\ell/2}\).  For \(\ell=1\), \(w_{\{i\}}=\alpha_i\),
so the entropy expression cancels exactly.
\end{proof}

\begin{theorem}[Grouped hubs with an arbitrary-rank crossing remainder]
\label{thm:grouped-crossing}
Let \(X_H\) and \(X'\) be independent, with \(X'\) having
independent coordinate densities bounded by one.
Let \(U\) have \(k\) orthonormal columns, and suppose
\(W=U^{\mathsf T}X_H\) has density at most \(M_H\).
Consider
\[
 Q=F(X_H)+2\langle W,MX'+\beta\rangle+R(X'),
\]
where \(F\) is arbitrary measurable, \(M\) has row rank
\(k\), and \(R\) is an arbitrary quadratic polynomial
with matrix \(A_R\).
Split the light coordinates into \(I,J\), assume
\(M_J\) has row rank \(k\), and write \(S\) for the
product of its singular values.  Define
\[
 C=(A_R)_{IJ}(I-\Pi_{\operatorname{row}(M_J)}).
\]
Suppose \(Z_\ell(C)>0\) and its \(\ell\)-volume row
marginals are at most \(1/2\).  Put \(d=k+\ell\).
For \(t\ne0\),
\[
 |\phi_Q(t)|\leq\min\{1,K|t|^{-d/2}\},
 \quad
 K=\frac{M_H^{1/2}2^{(k+3\ell)/4}\pi^{d/2}e^{\ell/4}}
 {S^{1/2}Z_\ell(C)^{1/4}}.
\]
If \(d>2\), then
\begin{equation}
 \|p_Q\|_\infty\leq
 \frac d{d-2}
 2^{(k+3\ell)/(2d)}e^{\ell/(2d)}
 \frac{M_H^{1/d}}{S^{1/d}Z_\ell(C)^{1/(2d)}}.
 \label{eq:grouped-crossing-density}
\end{equation}
For \(\ell=1\), omit the exponential factors.
\end{theorem}
\begin{proof}
Push \(e^{itF(X_H)}d\mathbb P_H\) onto \(W\), obtaining
a complex density \(a_t\) with
\(\|a_t\|_2^2\leq M_H\), and put
\(h_t(l)=\widehat a_t(2t(l+\beta))\).
Plancherel gives
\[
 \int|h_t(l)|^2\,dl\leq
 M_H(\pi/|t|)^k.
\]
Let \(\mu_t(l)\) be the density of the complex measure
\(\mathbb E[e^{itR(X')};MX'\in dl]\).
The ordinary projection density of \(MX'\) is bounded,
so Lemma~\ref{lem:phase-projection} implies \(\mu_t\in L^2\).
Independence yields
\[
 \phi_Q(t)=\int h_t(l)\mu_t(l)\,dl,\qquad
 |\phi_Q(t)|^2\leq M_H(\pi/|t|)^k\|\mu_t\|_2^2.
 \label{eq:grouped-pairing}
\]

To bound the last norm, take a nonnegative even compactly
supported approximate identity \(k_\varepsilon\) on
\(\mathbb R^k\) and write
\[
 m_\varepsilon(l)=
 \mathbb E[e^{itR(X')}k_\varepsilon(l-MX')].
\]
Conditional Cauchy--Schwarz in \(X_I\), followed by
integration in \(l\) and an independent copy \(X_J^*\),
gives
\[
 \|m_\varepsilon\|_2^2
 \leq
 \mathbb E\left[
 e^{it(R(X_I,X_J)-R(X_I,X_J^*))}
 \kappa_\varepsilon(M_JD_J)\right],
\]
where \(D_J=X_J-X_J^*\) and
\(\kappa_\varepsilon=k_\varepsilon*k_\varepsilon\)
is a nonnegative probability density.
Every term purely in \(I\) cancels.  Terms purely in
\(J\) give a phase of modulus one, while the cross terms
give \(2\langle X_I,(A_R)_{IJ}D_J\rangle\).
Averaging over \(X_I\) and taking absolute values bounds
the last expression by
\[
 \mathbb E\left[
 \kappa_\varepsilon(M_JD_J)
 \prod_{i\in I}
 |\widehat f_i(2t\langle(A_R)_{iJ},D_J\rangle)|
 \right].
\]
The coordinates of \(D_J\) are independent with densities
\(f_j*\widetilde f_j\leq1\).
Apply Lemma~\ref{lem:weighted-subset}, uniformly in
\(\varepsilon\), and pass to the limit by \(L^2\)
approximate-identity convergence.  This yields
\[
 \|\mu_t\|_2^2\leq
 \frac{2^{(k+\ell)/2}(2\pi)^\ell e^{\ell/2}}
 {S\sqrt{Z_\ell(C)}\,|t|^\ell}.
\]
Combining with the hub Plancherel estimate proves the
characteristic-function bound.  The integration formula
in Theorem~\ref{thm:exact-crossblock} gives
\eqref{eq:grouped-crossing-density}.
\end{proof}

\begin{corollary}[One hub and a spread second crossing volume]
\label{cor:one-hub-arbitrary-crossing}
For one original hub, let
\[
 Q=F(X_h)+2X_h(\langle w,X'\rangle+\beta)+R(X').
\]
For a split \(I,J\), suppose \(w_J\ne0\), and let
\(C=(A_R)_{IJ}(I-\Pi_{w_J})\).
If \(Z_2(C)>0\) and every second-volume row marginal is
at most \(1/2\), then
\[
 \|p_Q\|_\infty
 \leq3\,2^{7/6}e^{1/3}
       (|w_J|^2Z_2(C))^{-1/6}
 <9.4(|w_J|^2Z_2(C))^{-1/6}.
\]
There is no restriction on \(\operatorname{rank}C\).
\end{corollary}
\begin{proof}
Use Theorem~\ref{thm:grouped-crossing} with
\(k=1,\ell=2\), and \(M_H=1\).
\end{proof}

\begin{corollary}[Many coordinates with two exact hub directions]
\label{cor:many-two-hubs}
Suppose a hub--light cross block has exact rank two and
write it as \(UM\), with \(U^{\mathsf T}U=I_2\).
Every internal hub quadratic, and even every measurable
internal hub function, is allowed.
If a light split has \(\operatorname{rank}M_J=2\),
singular-value product \(S\), and projected crossing
matrix \(C\) with
\[
 F_0=\|C\|_{\mathrm F}>0,\qquad
 \|c_i\|^2\leq F_0^2/2,
\]
then product hub and light laws satisfy
\[
 \|p_Q\|_\infty
 \leq3\,2^{7/6}(SF_0)^{-1/3}<6.74(SF_0)^{-1/3}.
\]
\end{corollary}
\begin{proof}
The projection \(U^{\mathsf T}X_H\) has density at most
two.  Apply Theorem~\ref{thm:grouped-crossing} with
\(k=2,\ell=1\), \(M_H=2\), and \(Z_1(C)=F_0^2\).
\end{proof}

\section{Concentrated directions and uniform bounds}
\label{sec:structural-boundaries}

\begin{lemma}[Few concentrated directions in exact finite rank]
\label{lem:heavy-isotropic}
Let \(A=RMR^{\mathsf T}\), with \(R^{\mathsf T}R=I_r\)
and every eigenvalue of \(M\) of modulus at least one.
Assume \(A\) has no \(c\)-pivot.  Let \(H\) be any set of
original coordinates and
\[
 G_H=R_H^{\mathsf T}R_H,\qquad
 \tau=|H|\sqrt{2c}.
\]
If the eigenvalues of \(G_H\) are \(g_1\geq\cdots\geq g_r\),
then \(g_k\leq\tau\) whenever \(k>r/2\).
If \(M\) is definite, every row leverage
\(|R^{\mathsf T}e_i|^2\) is less than \(c\).
\end{lemma}
\begin{proof}
By \eqref{eq:no-pivot-entries},
\(\|A_{HH}\|\leq\|A_{HH}\|_{\mathrm F}\leq\tau\).
Rotate the frame so that \(G_H\) is diagonal.
If \(g_k>\tau\), the compression of \(M\) to its first
\(k\) frame coordinates has norm at most \(\tau/g_k<1\).
Indeed, writing \(R_H=VG_H^{1/2}\) on its support
expresses \(A_{HH}\) as an isometric copy of
\(G_H^{1/2}MG_H^{1/2}\), and division by the first
\(k\) eigenvalues of \(G_H^{1/2}\) gives this norm bound.
Interlacing then forces \(M\) to have at least \(k\)
eigenvalues greater than \(-1\), hence at least \(1\),
and at least \(k\) eigenvalues less than \(1\), hence at
most \(-1\).  These are disjoint sets, so \(2k\leq r\).
For definite \(M\), its eigenvalue hypothesis gives
\(|a_{ii}|\geq|R^{\mathsf T}e_i|^2\); use the
no-pivot diagonal bound.
\end{proof}

For example, if \(H=\{i:|R^{\mathsf T}e_i|^2\geq\delta\}\),
then \(|H|\leq r/\delta\).  In exact rank three,
Lemma~\ref{lem:heavy-isotropic} permits only one strongly
concentrated frame direction when \(c\) is small relative
to \(\delta\).  It makes no assertion about a decomposition
with an arbitrary additional diagonal: that diagonal may
cancel the diagonal entries of the low-rank component.

\begin{proposition}[A normalized exact-rank-three transition]
\label{prop:rank-three-epsilon}
Let \(w,v,p\) be orthonormal vectors on the light
coordinates, write
\(L=\langle w,X'\rangle\), \(V=\langle v,X'\rangle\),
\(P=\langle p,X'\rangle\), and put
\[
 Q_\varepsilon=2X_hL+V^2+2\varepsilon LP.
\]
For every real \(\varepsilon\), its quadratic matrix has
exact rank three and nonzero eigenvalues
\[
 \sqrt{1+\varepsilon^2},\qquad
 -\sqrt{1+\varepsilon^2},\qquad 1.
\]
For every fixed \(c>0\) and \(0\leq\varepsilon\leq1\),
there are arbitrarily large spread frames for which this
matrix has no \(c\)-pivot.
If the light frame entries all have modulus \(m^{-1/2}\),
then every principal submatrix of order at most \(N\)
has third singular value at most \(N/m\).
\end{proposition}
\begin{proof}
In the orthonormal basis \((e_h,w,v,p)\), the matrix is
\[
 \begin{pmatrix}
 0&1&0&0\\
 1&0&0&\varepsilon\\
 0&0&1&0\\
 0&\varepsilon&0&0
 \end{pmatrix}.
\]
On the plane spanned by \(w\) and
\((e_h+\varepsilon p)/\sqrt{1+\varepsilon^2}\), it is
\(\bigl(\begin{smallmatrix}0&\sqrt{1+\varepsilon^2}\\
\sqrt{1+\varepsilon^2}&0\end{smallmatrix}\bigr)\).
The \(v\)-direction has eigenvalue one and the remaining
orthogonal direction has eigenvalue zero.

Take three orthogonal Walsh columns on \(m\) coordinates,
normalized to entries of modulus \(m^{-1/2}\), with
\(m>6/c\).  The hub diagonal is zero, every hub--leaf
minor has modulus \(1/m<c\), and light-block entries
have modulus at most \(3/m\).
Their diagonals are less than \(c\) and their principal
two-coordinate determinants have modulus at most
\(18/m^2<c\), after enlarging \(m\) if necessary.
This proves the no-pivot assertion.

For any coordinate set \(K\), the principal matrix is
the sum of a matrix of rank at most two,
\[
 \bigl((e_h+\varepsilon p)w^{\mathsf T}
       +w(e_h+\varepsilon p)^{\mathsf T}\bigr)_{KK},
\]
and \(v_Kv_K^{\mathsf T}\), of norm
\(|v_K|^2\leq |K|/m\).
The singular-value approximation inequality gives
\(s_3(A_{KK})\leq |K|/m\).
\end{proof}

The family in Proposition~\ref{prop:rank-three-epsilon}
has a third singular value of one for every \(\varepsilon\).
The endpoint \(\varepsilon=0\) and every fixed nonzero
\(\varepsilon\) are covered by results in
Section~\ref{sec:hubs}, but the available bounds for
arbitrary rough laws need not stay bounded as
\(\varepsilon\) tends to zero.
Corollary~\ref{cor:hub-loglog-density} gives a much smaller
loss than the direct leakage estimate.
Theorem~\ref{thm:grouped-crossing} also retains a dependence
on this parameter: in the one-hub split its projected
crossing matrix is
\[
 C=v_I(v_J^\perp)^{\mathsf T}
       +\varepsilon w_I(p_J^\perp)^{\mathsf T},
\]
so \(Z_2(C)\leq4\varepsilon^2\) for
\(0\leq\varepsilon\leq1\).
The finite-core theorem is likewise insufficient uniformly,
by the last assertion of Proposition~\ref{prop:rank-three-epsilon}.
These are limitations of the respective estimates, not
counterexamples to a bounded density.

\begin{proposition}[Two hubs do not suffice with an arbitrary remainder]
\label{prop:two-hub-counterexample}
For every \(0<c\leq1\), there is a real symmetric matrix
with no \(c\)-pivot, exactly two eigenvalues of modulus
at least one, and a two-row hub--light block whose
singular values are both one, such that the corresponding
quadratic form under independent standard Gaussian
coordinates has unbounded density.
\end{proposition}
\begin{proof}
Choose \(0<b<c/2\), put
\(B=\operatorname{diag}(b,-b)\), and let \(M\) have
two orthonormal, sufficiently spread rows.
Set
\[
 A=\begin{pmatrix}
 B&M\\ M^{\mathsf T}&M^{\mathsf T}B^{-1}M
 \end{pmatrix}.
\]
Writing \(L=MX'\),
\[
 Q_A=(X_H+B^{-1}L)^{\mathsf T}
       B(X_H+B^{-1}L).
\]
The matrix has rank two.  On the two mutually orthogonal
hub--partner planes its nonzero eigenvalues are
\(b+b^{-1}\) and \(-b-b^{-1}\).
By taking the two rows of \(M\) with entries of modulus
at most \(\eta\), where
\(\eta<c/2\) and \(2\eta^2/b<c/2\), every entry of
\(A\) has modulus less than \(c/2\).  Thus every
diagonal entry is smaller than \(c\) and every principal
two-coordinate determinant has modulus less than
\(c^2/2<c\).

For Gaussian inputs, \(Y=X_H+B^{-1}L\) is a nondegenerate
Gaussian vector whose density is bounded below by some
\(\gamma>0\) on a neighbourhood of the origin.
Use the change of variables
\(u=\sqrt b(Y_1-Y_2)\),
\(v=\sqrt b(Y_1+Y_2)\).
Then \(Q_A=uv\) and the Jacobian is \(2b\).
Choose \(r>0\) so that the square \(|u|,|v|\leq r\)
maps into that neighbourhood.  For \(0<a\leq r^2\),
\[
 \mathbb P\{|Q_A|\leq a\}
 \geq\frac{\gamma}{2b}
 \left|\{(u,v):|u|,|v|\leq r,\ |uv|\leq a\}\right|
 =\frac{2\gamma}{b}a(1+\log(r^2/a)).
\]
The ratio to \(2a\) is unbounded as \(a\downarrow0\).
\end{proof}

Proposition~\ref{prop:two-hub-counterexample} has only two
large eigenvalues.  It shows why a two-hub bound must use
additional information about the remainder, such as that
in Theorem~\ref{thm:grouped-crossing}; it does not refute
the three-eigenvalue conjecture.

The structural statements above do not imply that every
no-pivot matrix admits a split for which the quantities in
\eqref{eq:grouped-crossing-density} are uniformly nondegenerate.
In particular, the grouped hub block must have exact rank;
discarding the smaller singular directions does not cover a
nearly rank-two block.  A dimension-independent bound requires
uniform control of these limiting configurations and of the
one-hub transition above.
